\documentclass[10pt]{beamer}
\usetheme{Madrid}
\usepackage{booktabs}
\usepackage{graphicx}
\usepackage[T1]{fontenc}
\setbeamertemplate{navigation symbols}{}

% ======================= EDIT YOUR DETAILS =======================
\newcommand{\studentname}{Aflah Muhammed P}
% Change the university roll/registration number:
\newcommand{\rollno}{TVE23CSXXX}
% Change your class roll number:
\newcommand{\rollclass}{Roll No: XX}
% Change your guide name:
\newcommand{\guidename}{Prof. Guide Name}
% Change department / college / short-institute if needed:
\newcommand{\deptname}{Department of Computer Science and Engineering}
\newcommand{\collegename}{College of Engineering Trivandrum}
\newcommand{\shortinst}{CET}
% Change the seminar date:
\newcommand{\seminardate}{\today}
% =================================================================

\title[B.Tech Seminar]{Large Language Models for Planning: How AI Agents Turn a Goal into Step-by-Step Actions}
\subtitle{Based on: Large Language Models for Planning: A Comprehensive and Systematic Survey (arXiv:2505.19683, 2025)}
\author[\studentname\ (\shortinst)]{\studentname \\ \small \rollno \\ \small \rollclass \\ \small Guide: \guidename}
\institute[]{\deptname \\ \collegename}
\date{\seminardate}

\begin{document}

\begin{frame}[plain]
  \titlepage
\end{frame}

\begin{frame}{Table of Contents}
  \tableofcontents
\end{frame}

\section{Introduction}
\begin{frame}{Introduction}
\begin{itemize}
  \item Planning means turning a high-level goal into ordered, executable actions
  \item Intelligent agents need planning to complete real multi-step tasks
  \item LLMs can plan from plain-language goals without hand-written rules
  \item But LLMs hallucinate steps and struggle on long tasks
  \item This survey maps the fast-growing field of LLM-based planning
\end{itemize}
\end{frame}

\section{Literature Survey}
\begin{frame}{Literature Survey / Background}
  \small
  \begin{tabular}{@{}p{2.7cm}p{2.3cm}p{5.6cm}@{}}
    \toprule
    \textbf{Title} & \textbf{Author} & \textbf{Summary} \\
    \midrule
    Chain-of-Thought Prompting Elicits Reasoning in LLMs & Wei et al., 2022 & Shows that prompting a model to reason step by step greatly improves multi-step reasoning, a foundation for LLM planning. \\
    \addlinespace
    ReAct: Synergizing Reasoning and Acting in Language Models & Yao et al., 2023 & Interleaves reasoning with actions so an LLM agent can plan and interact with an environment at the same time. \\
    \addlinespace
    Tree of Thoughts: Deliberate Problem Solving with LLMs & Yao et al., 2023 & Lets an LLM explore and evaluate multiple reasoning branches like a search tree instead of one linear guess. \\
    \addlinespace
    \bottomrule
  \end{tabular}
\end{frame}

\section{Motivation}
\begin{frame}{Motivation}
\begin{itemize}
  \item LLM agents are moving into web, coding, and robotics tasks
  \item Hundreds of planning methods exist with no shared map
  \item Each method makes different trade-offs and assumptions
  \item Beginners and researchers need one organized, systematic overview
\end{itemize}
\end{frame}

\section{Problem Statement}
\begin{frame}{Problem Statement}
  \begin{block}{}
  LLMs can generate plans from natural-language goals but often produce invalid, inefficient, or incomplete action sequences, and the many proposed fixes lack a unified organization and evaluation view. This survey systematically categorizes LLM-based planning methods and their benchmarks and metrics.
  \end{block}
\end{frame}

\section{Objectives}
\begin{frame}{Objectives}
\begin{itemize}
  \item Establish clear definitions and categories of automated planning
  \item Provide a systematic taxonomy of LLM-based planning methods
  \item Organize methods into three principal approach families
  \item Summarize benchmarks, metrics, and performance comparisons
  \item Discuss enabling mechanisms and future research directions
\end{itemize}
\end{frame}

\section{Methodology}
\begin{frame}[allowframebreaks]{Methodology / Approach}
\begin{itemize}
  \item Survey and taxonomy paper, not a single new algorithm
  \item Start with foundations of classical automated planning
  \item Family 1: External Module Augmented, add planners or tools
  \item Family 2: Finetuning-based, train on trajectories and feedback
  \item Family 3: Searching-based, decompose tasks and search plan space
  \item Consolidate evaluation benchmarks, metrics, and open challenges
\end{itemize}
\end{frame}

\section{Key Concepts}
\begin{frame}[allowframebreaks]{Key Concepts}
  \begin{description}
    \item[External Module Augmented] Pair the LLM with outside helpers like classical planners, tools, or memory to improve plans.
    \item[Finetuning-based Methods] Adjust the LLM weights using planning trajectories and success or failure feedback signals.
    \item[Searching-based Methods] Decompose complex tasks and explore multiple candidate plans via search or better decoding.
    \item[Task decomposition] Split a hard goal into smaller subtasks that are easier to solve step by step.
    \item[Evaluation framework] Shared benchmarks and metrics used to compare planning methods fairly.
  \end{description}
\end{frame}

\section{System Architecture}
\begin{frame}{System Architecture / How It Works}
  \begin{block}{Overview}
  A diagram can show a single natural-language goal entering an LLM planner at the top. From the LLM, three parallel branches represent the three families: one branch sends structured output to an external classical planner or tool and returns a validated plan; a second branch shows the LLM being fine-tuned offline on trajectory and feedback data before it plans; a third branch shows the LLM decomposing the goal into subtasks and expanding a search tree of candidate steps that are scored and pruned. All three branches converge on a produced action sequence, which is then executed in an environment and scored by an evaluation block measuring success rate, executability, and efficiency.
  \end{block}
  % TIP: insert a diagram here, e.g.:
  % \begin{center}\includegraphics[width=0.7\textwidth]{your-figure.png}\end{center}
\end{frame}

\section{Results}
\begin{frame}[allowframebreaks]{Results / Findings}
\begin{itemize}
  \item Consolidates LLM planning research into three clear families
  \item External modules add guarantees but need formal task descriptions
  \item Fine-tuning bakes in planning skill but needs trajectory data
  \item Search improves plan quality at extra computation cost
  \item Reviews benchmarks like Blocksworld, ALFWorld, WebShop, TravelPlanner
\end{itemize}
\end{frame}

\section{Discussion}
\begin{frame}{Discussion}
\begin{itemize}
  \item No single family wins; each trades off cost, data, and guarantees
  \item Combining families often yields the strongest agents
  \item Reliable planning is key to safe, useful autonomous agents
  \item Shared benchmarks and metrics enable fair method comparison
\end{itemize}
\end{frame}

\section{Challenges}
\begin{frame}{Limitations \& Challenges}
\begin{itemize}
  \item LLM planners still hallucinate invalid or unsafe steps
  \item Performance drops on long-horizon and highly complex tasks
  \item Benchmarks may not reflect messy real-world environments
  \item Search and fine-tuning add computational or data cost
\end{itemize}
\end{frame}

\section{Future Work}
\begin{frame}{Future Work}
\begin{itemize}
  \item Scale planning to longer horizons and harder domains
  \item Improve robustness when environments shift unexpectedly
  \item Build richer, more realistic evaluation benchmarks
  \item Blend classical planning guarantees with neural flexibility
\end{itemize}
\end{frame}

\section{Conclusion}
\begin{frame}{Conclusion}
\begin{itemize}
  \item LLMs make planning accessible directly from plain-language goals
  \item The field fits into three families: external modules, fine-tuning, search
  \item Shared benchmarks and metrics let us compare methods fairly
  \item Reliable LLM planning is central to future capable, safe agents
\end{itemize}
\end{frame}

\section{References}
\begin{frame}[allowframebreaks]{References}
  \footnotesize
  \begin{enumerate}
    \item Large Language Models for Planning: A Comprehensive and Systematic Survey, Cao P., Men T., Liu W., Zhang J., Li X., Lin X., Sui D., Cao Y., Liu K., Zhao J., arXiv:2505.19683, 2025
    \item ReAct: Synergizing Reasoning and Acting in Language Models, Yao S. et al., ICLR, 2023
    \item Tree of Thoughts: Deliberate Problem Solving with Large Language Models, Yao S. et al., NeurIPS, 2023
    \item Chain-of-Thought Prompting Elicits Reasoning in Large Language Models, Wei J. et al., NeurIPS, 2022
    \item LLM+P: Empowering Large Language Models with Optimal Planning Proficiency, Liu B. et al., arXiv:2304.11477, 2023
    \item Reflexion: Language Agents with Verbal Reinforcement Learning, Shinn N. et al., NeurIPS, 2023
  \end{enumerate}
\end{frame}

\begin{frame}[plain]
  \begin{center}
    {\Huge Thank You}\\[1em]
    {\large Questions?}
  \end{center}
\end{frame}

\end{document}